\begin{table*}[!t]
\centering
\caption{Optional memory module (Mem), off by default and not used for the results of this paper: a phenomenological mark written where PIP3 is high and erased slowly. A coloured reaction is the one to which the equally coloured statement in the text and the equally coloured passage in the cited paper refer; a small square appends a further statement. $H(\mathrm{PIP3})=N_\mathrm{PIP3}^3/(K_\mathrm{mem}^3+N_\mathrm{PIP3}^3)$.}
\label{tab:mmemrxn}
\footnotesize
\renewcommand{\arraystretch}{1.15}
\begin{tabularx}{\textwidth}{@{}l>{\raggedright\arraybackslash}p{0.37\textwidth}>{\raggedright\arraybackslash}X@{}}
\toprule
ID & Reaction & Propensity (rate law)\\
\midrule
\mkt{Mmem-mem-carrier}{M.1} & \mkt{Mmem-mem-carrier}{$\mathrm{Mem}_c \to \mathrm{Mem}_m^*$} & $k_\mathrm{w}\,[\mathrm{Mem}_c]\,H(\mathrm{PIP3})$ \\
\mkn{Mmem-mem-dir}{Mmem-mem-carrier}{M.2} & \mkn{Mmem-mem-dir}{Mmem-mem-carrier}{$\mathrm{Mem}_m^* \to \mathrm{Mem}_c$} & $k_\mathrm{e}[\mathrm{Mem}_m^*]$ \\
\bottomrule
\end{tabularx}
\end{table*}
